\subsection{Task B. Dự đoán thao tác chỉnh sửa đồ thị}
Cho đồ thị cảnh nguồn $G_s$ và câu yêu cầu chỉnh sửa $q$, mô hình dự đoán tập
thao tác $\Delta_q$ có thể áp dụng lên $G_s$. Mỗi thao tác chỉ rõ đối tượng
hoặc cạnh cần thêm, xóa hay thay. Với cạnh, mô hình phải xác định cả đỉnh đầu,
quan hệ và đỉnh cuối. Đổi màu của một vật được biểu diễn bằng thao tác thay đỉnh,
giữ mã đối tượng và cập nhật thuộc tính màu. Nhóm đặt hai câu hỏi thực nghiệm.
Trên MagicBrush, học từ những thay đổi phục vụ câu lệnh có cải thiện dự đoán
so với học từ toàn bộ chênh lệch giữa hai đồ thị hay không? Khi có sẵn cặp đồ
thị và câu lệnh, mô hình học được những thao tác nào trên bộ EMNLP 2020?

\subsubsection{Xây bộ dữ liệu từ MagicBrush}
MagicBrush~\cite{zhang2023magicbrush} cung cấp ảnh trước và sau chỉnh sửa
cùng câu lệnh. Nhóm dùng gán nhãn có hỗ trợ AI để mô tả đối tượng, thuộc tính
và quan hệ trong hai ảnh, sau đó ghép các thực thể tương ứng. Từ 437 cặp có
đồ thị, nhóm lập hai bộ nhãn cho cùng đầu vào. Nhãn \textit{chênh lệch đồ thị}
ghi mọi thao tác cần để tái tạo đồ thị của ảnh đích. Nhãn \textit{theo câu
lệnh} chỉ giữ các thay đổi phục vụ yêu cầu, kể cả thao tác cấu trúc cần thiết
như xóa cạnh nối với một đỉnh bị xóa. Khi hai nhãn khác nhau, ảnh đích là
bằng chứng để xem lại đối tượng và quan hệ. Lúc suy luận, mô hình chỉ nhận
$G_s$ và $q$.

\begin{figure}[!htbp]\centering
\begin{tabular}{@{}cc@{}}
\includegraphics[width=.455\linewidth]{magicbrush-laptop-source.jpg} &
\includegraphics[width=.455\linewidth]{magicbrush-laptop-target.jpg}\\
(a) Ảnh nguồn & (b) Ảnh sau chỉnh sửa
\end{tabular}\par\smallskip
\includegraphics[width=.98\linewidth]{magicbrush-laptop-graph.pdf}
\caption{Mẫu validation \texttt{484847:1} của MagicBrush với yêu cầu
\textit{put a laptop on the table}. Đồ thị của (a) có bàn, sách và điều khiển.
Để tái tạo đồ thị của (b) cần sáu thao tác. Yêu cầu chỉ xác định hai thao tác:
thêm đỉnh laptop và thêm cạnh $(\mathrm{laptop},\mathrm{on},\mathrm{table})$.
Sơ đồ bên dưới chỉ vẽ phần đồ thị liên quan đến bàn. Hai cạnh có sẵn nối sách và điều khiển với bàn được giữ lại.}
\label{fig:mb-intent-laptop}
\end{figure}

Nhóm đối chiếu từng cặp với câu lệnh như Hình~\ref{fig:mb-intent-laptop}.
Tổng cộng 92 cặp được để riêng do còn mơ hồ về đối tượng, quan hệ hoặc ảnh
đích chưa thể hiện đủ yêu cầu. Bộ cuối gồm \textbf{345 cặp} $(G_s,q,\Delta_q)$
(Bảng~\ref{tab:mb-cohort}), vượt mức tối thiểu 300 cặp của đề. Các lượt chỉnh
sửa thuộc cùng nhóm ảnh nguồn nằm trong cùng một tập. Nhóm cũng kiểm tra
hai vai ảnh trước/sau bằng mã băm pixel để phát hiện ảnh trùng giữa các tập.

\begin{table}[!htbp]\centering\small
\caption{Số cặp MagicBrush trước và sau khi đối chiếu theo câu lệnh.
Cột cuối đếm các cặp được giữ có nhãn khác chênh lệch đồ thị.}
\label{tab:mb-cohort}
\begin{tabular}{@{}lrrrr@{}}\toprule
Tập & Ban đầu & Giữ & Để riêng & Nhãn khác\\\midrule
Train & 301 & 240 & 61 & 205\\
Validation & 67 & 46 & 21 & 43\\
Test & 69 & 59 & 10 & 48\\\midrule
Tổng & 437 & 345 & 92 & 296\\\bottomrule
\end{tabular}\end{table}

\begin{figure*}[!t]\centering
\includegraphics[width=.93\textwidth]{magicbrush-test-errors.pdf}
\caption{Hai mẫu MagicBrush test, dự đoán của GNN học nhãn theo câu lệnh,
seed 42. Ảnh đích dùng để đối chiếu nhãn. Ở (a), mẫu \texttt{181753:1},
bàn giữ chất liệu gỗ sẫm và đổi từ chữ nhật sang tròn. Mô hình giữ nguyên đồ thị.
Ở (b), mẫu \texttt{65564:2}, bóng bay xuất hiện phía trên hươu. Câu lệnh chỉ
yêu cầu thêm bóng bay. Mô hình thêm vịt và nối nó với đỉnh hươu ở tiền cảnh.}
\label{fig:mb-test-errors}
\end{figure*}

Nhóm lấy ngẫu nhiên 20 cặp từ 345 cặp, không hoàn lại (seed 20261119).
Người duyệt xem hai ảnh, câu lệnh, đồ thị nguồn và toàn bộ thao tác đích.
Một cặp đạt khi các thao tác tác động đúng thực thể, thực hiện đủ yêu cầu
và bảo toàn nội dung ngoài yêu cầu. Người duyệt chấp nhận \textbf{20/20 cặp
(100\%)}. Tỷ lệ này mô tả mức chấp nhận của một người trên mẫu kiểm tra tay.
Cả 345 nhãn cũng vượt kiểm tra tự động về khả năng áp dụng thao tác lên đồ
thị nguồn, bao gồm kiểm tra hai đầu mút của mỗi cạnh.

\subsubsection{Huấn luyện và đối chứng hai cách gán nhãn}
Nhóm dùng GNN hai lớp, chiều ẩn 192. CLIP text ViT-B/32 đóng băng mã hóa câu
lệnh, nhãn và thuộc tính đỉnh, cùng quan hệ trên cạnh. Bộ dự đoán chọn thao
tác trên thực thể nguồn hoặc tạo thực thể mới, chọn nhãn và thuộc tính đỉnh,
rồi chọn quan hệ và hai đầu mút cho cạnh. Hai nhánh dùng cùng 240 cặp train,
chung từ vựng xây từ train, cùng khởi tạo và thứ tự batch theo từng seed
42, 123, 3407. Biến thay đổi giữa hai nhánh là nhãn huấn luyện.
Nhóm dùng AdamW, tốc độ học $3\times10^{-4}$, batch 16. Cả hai nhánh chọn
checkpoint theo micro-F1 thao tác trên 46 cặp validation có nhãn theo câu lệnh.
Bộ chấm tính một thao tác đúng khi loại, thực thể và nội dung chỉnh sửa đều
khớp sau khi chuẩn hóa mã thực thể mới. Trong thử nghiệm ban đầu trên 69
cặp test dùng toàn bộ chênh lệch làm nhãn chuẩn, MLP và GNN lần lượt đạt
micro-F1 3,75\% và 7,34\%. Thực nghiệm dưới đây dùng cùng 59 cặp và
cùng nhãn chuẩn theo câu lệnh để so sánh hai cách huấn luyện.

\begin{table}[!htbp]\centering\small
\caption{MagicBrush test. Precision (P), recall (R) và F1 theo phần trăm,
trung bình ba seed. Cả hai nhánh cùng được chấm theo câu lệnh trên 59 cặp.
$N$ là số thao tác chuẩn trên tập test.}
\label{tab:mb-prf}
\setlength{\tabcolsep}{3pt}
\begin{tabular}{@{}lr rrr rrr@{}}\toprule
& & \multicolumn{3}{c}{Chênh lệch đồ thị} & \multicolumn{3}{c}{Theo câu lệnh}\\
\cmidrule(lr){3-5}\cmidrule(l){6-8}
Thao tác & $N$ & P & R & F1 & P & R & F1\\\midrule
Thêm đỉnh & 24 & 0,00 & 0,00 & 0,00 & 5,95 & 5,56 & 5,73\\
Xóa đỉnh & 12 & 26,51 & 16,67 & 20,29 & 16,67 & 2,78 & 4,76\\
Thay đỉnh & 35 & 0,00 & 0,00 & 0,00 & 0,00 & 0,00 & 0,00\\
Thêm cạnh & 17 & 2,39 & 5,88 & 3,39 & 2,56 & 1,96 & 2,22\\
Xóa cạnh & 14 & 25,50 & 40,48 & 30,68 & 35,68 & 14,29 & 18,84\\
Thay cạnh & 0 & \nd & \nd & \nd & \nd & \nd & \nd\\\midrule
Micro & 102 & 5,30 & 8,50 & 6,53 & 5,35 & 3,92 & 4,45\\\bottomrule
\end{tabular}\end{table}

Trên 59 cặp test, micro-F1 của nhánh chênh lệch đồ thị đạt
$6,53\pm1,75\%$, nhánh theo câu lệnh đạt $4,45\pm1,59\%$
(trung bình $\pm$ độ lệch chuẩn mẫu qua ba seed). Nhánh theo câu lệnh có
micro-F1 thấp hơn ở cả ba seed. Precision trung bình gần bằng nhau, nhưng
recall giảm từ 8,50\% xuống 3,92\%. Bảng~\ref{tab:mb-prf} cho thấy chiều
thay đổi tùy thao tác. Học theo câu lệnh giúp F1 thêm đỉnh tăng lên
5,73\%, trong khi F1 xóa cạnh giảm từ 30,68\% xuống 18,84\%.
Cả hai nhánh đều có F1 bằng 0 cho 35 thao tác thay đỉnh. Riêng thao tác
thay cạnh có ba nhãn ở train và không có nhãn ở test. Dấu gạch trong bảng
biểu thị thiếu mẫu để đánh giá.

Toàn bộ tập thao tác dự đoán đều áp dụng được về mặt cấu trúc. Tỷ lệ khớp
chính xác đồ thị đích theo câu lệnh chỉ đạt 0,00\% với nhãn chênh lệch và
0,56\% với nhãn theo câu lệnh, trung bình ba seed. Nếu quy ước yêu cầu
đơn giản là cặp có đúng một thao tác đích, cả hai nhánh đạt 0/34 mẫu
khớp toàn bộ tập thao tác ở từng seed. Hình~\ref{fig:mb-test-errors} minh họa hai lỗi khác nhau của cùng
một checkpoint: bỏ sót thay đổi thuộc tính và chọn sai tên đối tượng khi thêm đỉnh.

\subsubsection{Hướng mở rộng trên EMNLP 2020}
Để khảo sát khả năng học thao tác khi có sẵn đồ thị nguồn, câu lệnh và đồ
thị đích, nhóm huấn luyện riêng MLP và GNN trên bộ chỉnh sửa đồ thị của
He và cộng sự~\cite{he2020sgmod}. Bộ này có đồ thị nguồn nhỏ (1--3 đỉnh).
Nhóm dùng 29.868 mẫu train, 532 validation và 530 test. Từ 1.000 mẫu test
chính thức, nhóm loại 468 mẫu trùng đồ thị nguồn với train hoặc validation,
cùng hai mẫu có ghép cặp thực thể mơ hồ. Việc chia tập dựa trên đồ thị nguồn;
metadata hiện có chưa đủ để kiểm chứng tính độc lập theo ảnh gốc.
Mỗi mô hình chạy với ba seed 42, 43, 44 và chọn checkpoint theo loss trên
validation.

\begin{table}[!htbp]\centering\small
\caption{EMNLP 2020, trung bình ba seed trên 530 mẫu test (\%). Macro-F1
lấy trung bình đều trên sáu loại thao tác. Cột cuối yêu cầu toàn bộ đồ thị
đích khớp nhãn.}\label{tab:emnlp}
\begin{tabular}{@{}lrrr@{}}\toprule
Mô hình & Micro-F1 & Macro-F1 & Khớp đồ thị đích\\\midrule
MLP & $81,69\pm0,31$ & $54,25\pm0,51$ & $67,74\pm0,68$\\
GNN & $82,44\pm0,39$ & $55,03\pm0,28$ & $68,55\pm0,22$\\\bottomrule
\end{tabular}\end{table}

\begin{figure}[!htbp]\centering
\includegraphics[width=.98\linewidth]{emnlp-gnn-operations.pdf}
\caption{F1 của GNN trên EMNLP theo từng thao tác, trung bình ba seed.
Số trong ngoặc là số thao tác chuẩn trên tập test. Phép xóa có nhiều mẫu
và đạt F1 cao ở cả đỉnh lẫn cạnh.}
\label{fig:emnlp-gnn-operations}
\end{figure}

Trong Hình~\ref{fig:emnlp-gnn-operations}, GNN đạt F1 97,76\% cho xóa đỉnh và 97,64\% cho xóa cạnh.
Hai loại này chiếm 1.091 trong 1.533 thao tác chuẩn, nên ảnh hưởng lớn đến
micro-F1 ở Bảng~\ref{tab:emnlp}. Thay đỉnh chỉ đạt 27,38\% trên 121 thao
tác. Cả sáu thao tác thay cạnh trong test đều bị dự đoán sai.
Kết quả cho thấy mô hình học tốt phép xóa trong thiết lập có đồ thị cho sẵn.
Quy mô train, kích thước đồ thị, phân bố thao tác và kiến trúc dự đoán
đều khác giữa hai bộ. Vì vậy, khoảng cách điểm giữa EMNLP và MagicBrush
chưa cho phép định lượng riêng ảnh hưởng của nhiễu nhãn.
Phép đối chứng trực tiếp trên MagicBrush cho thấy với 240 cặp train, hiệu
chỉnh nhãn để bám sát câu lệnh vẫn chưa đủ để cải thiện micro-F1.
\FloatBarrier
